\section{Settings and terminology}
\label{intro:sec:settings}

With these motivations in place, we now describe the concrete setting that serves as
our laboratory throughout the thesis.
We work in \emph{supervised learning with scalar labels}: given a dataset of
input--output pairs, the goal is to learn a function that predicts the label from
the input.
Among all learning paradigms, supervised regression is the one that affords the most
control: we do not need to rely on complex internal structure of the data for building our model,
and having a scalar label carries more signal compared to simple classification.
This level of control is precisely what a physicist needs to build and test a theory.

For any supervised learning algorithm we need three ingredients: a
\emph{data distribution}, a \emph{model class}, and a \emph{learning algorithm}.
Each ingredient offers enormous freedom, and our choices are guided by a single
criterion: retain only the features that are essential for the phenomena we wish to
study, and strip everything else.

\paragraph{What data do we need?}
Real-world data---images, text, audio---has rich statistical structure that is not
yet amenable to exact analysis.
We therefore replace it with the simplest distribution that still supports
non-trivial learning: \emph{isotropic Gaussian} inputs.
Concretely, $\vec{x} \in \mathbb{R}^d$ is drawn i.i.d.\ from
\[\vec{x} \sim \mathcal{N}(0, I_d).\]
Obviously, this is a strong assumption that does not reflect the complexity of real data.
The main analytical payoff is that any linear transformation of a Gaussian vector is
again Gaussian. As we will see, this turns a $d$-dimensional stochastic problem into a
low-dimensional one. Although real data is rarely Gaussian, Gaussian universality justifies it as a tractable proxy for studying SGD dynamics. \citet{goldt2020modeling}, made rigorous by \citet{goldt2022gaussian}, shows via the Gaussian Equivalence Property that for data drawn from a one-layer generative model (the hidden-manifold model) the pre-activations driving online SGD become jointly Gaussian in the high-dimensional limit, with the data entering only through its first two moments; the resulting learning dynamics then coincide with that on an equivalent Gaussian model. We therefore adopt i.i.d. Gaussian inputs, not as a literal data model but as the canonical representative of a universality class capturing the relevant high-dimensional dynamics.
There have been some recent efforts to relax this assumption: \citet{zweig2023single} generalize some single-index results to spherically symmetric distributions, while \citet{cagnetta2024how} introduce a different hierarchical model of the data. In independent component analysis the departure from Gaussianity is not a nuisance but the very quantity to be recovered---the hidden directions surface only through the higher-order cumulants that a Gaussian lacks~\citep{ricci2025feature}. The Gaussian case nonetheless remains the most tractable and best understood.

We still need a label $y$ for each input $\vec{x}$.
We take it to be generated by a fixed \emph{teacher} function $f^\star$, with
optional additive Gaussian noise:
\begin{equation}
\label{eq:intro:teacher}
    y = f^\star(\vec{x}) + \sqrt{\Delta}\,\xi, \qquad \xi \sim \mathcal{N}(0,1),
\end{equation}
where $\Delta \geq 0$ is the noise variance and $\Delta = 0$ is the \emph{realizable}
case; we add this to simulate the noise and imperfections of real data.
The teacher models the hidden rule that generates the data.
We would like it to be non-trivial but not arbitrary; the right compromise is to
suppose that the labels depend on the input only through a few
directions.
More precisely, we take the teacher to be a \emph{multi-index model}:
\begin{equation}
\label{eq:intro:multi-index}
    f^\star(\vec{x}) = h^\star\!\left(\frac{W^\star \vec{x}}{\sqrt{d}}\right),
\end{equation}
where $W^\star = \{\vec{w}^\star_r\}_{r \in [k]} \in \mathbb{R}^{k \times d}$ is the
\emph{teacher weight matrix} whose $k<d$ rows $\vec{w}^\star_r$ are the
\emph{teacher directions}, and $h^\star : \mathbb{R}^k \to \mathbb{R}$ is a possibly
nonlinear \emph{link function}.
\addnote{The normalization of the teacher directions plays an important role in the definition of all the subsequent quantities, and different conventions are used in the literature. We take them to be standard Gaussian, $\vec{w}^\star_r \sim \mathcal{N}(0, I_d)$, so that $\lVert\vec{w}^\star_r\rVert \approx \sqrt{d}$, and state this convention explicitly whenever it is relevant.}
The special case $k = 1$ is the \emph{single-index model}, with a single relevant
direction $\vec{w}^\star$.
The subspace $V^\star = \operatorname{span}(\vec{w}^\star_1, \ldots, \vec{w}^\star_k)$ in a
$d$-dimensional space is the \emph{teacher subspace}, and the learning process can be thought of as the discovery of this subspace and aligning with it.

\paragraph{What model do we use?}
Deep learning is a broad family of models, often designed towards a specific data modality or task.
Yet, the core of deep learning is the same across all these variations: stacked layers of linear transformations and nonlinear activations.
We take its simplest non-trivial representative: a \emph{two-layer (shallow) neural
network} with $p$ hidden neurons, first-layer weights
$W = \{\vec{w}_j\}_{j \in [p]} \in \mathbb{R}^{p \times d}$, second-layer weights
$\vec{a} = (a_1, \ldots, a_p)^\top \in \mathbb{R}^p$, and pointwise activation
function $\sigma : \mathbb{R} \to \mathbb{R}$:
\begin{equation}
\label{eq:intro:student}
    f(\vec{x};\, W, \vec{a})
    = \frac{1}{p} \sum_{j=1}^p a_j\,\sigma\!\left(\frac{\langle \vec{w}_j,
    \vec{x}\rangle}{\sqrt{d}}\right).
\end{equation}
Figure~\ref{intro:fig:teacher-student} contrasts this student network with the
multi-index teacher it is trained to imitate.

\begin{figure}[t]
  \centering
  \begin{minipage}[b]{0.49\linewidth}
    \centering
    \resizebox{\linewidth}{!}{% Schematic of the MULTI-INDEX teacher
%   f^*(x) = h^*(W^* x),   W^* in R^{k x d},  k < d  -- panel (a) of the
% teacher-student figure in chapters/introduction.tex.  Standalone tikzpicture,
% resized to its minipage at the call site.  The input is projected onto k
% directions (linear, no activation) and then passed jointly through a general
% nonlinear link h^*.  Needs tikz library arrows.meta (loaded in thesis.tex).
\begin{tikzpicture}[
    >={Latex[length=1.8mm]},
    font=\small,
    xnode/.style={circle, draw=netin!75, fill=netin!12, minimum size=7mm, inner sep=0pt},
    dnode/.style={circle, draw=#1!90, fill=#1!20, minimum size=8.5mm, inner sep=0pt},
    dnode/.default=nethid,
    hbox/.style={rounded corners=3pt, draw=netout!95, fill=netout!22,
                 minimum width=11mm, minimum height=21mm, inner sep=2pt},
    dots/.style={draw=none, fill=none, inner sep=1pt},
    wedge/.style={draw=#1!55, line width=0.5pt},
    wedge/.default=nethid,
    aedge/.style={draw=netout!75, line width=0.8pt},
    oedge/.style={->, draw=netout, line width=0.9pt},
    lab/.style={font=\footnotesize\itshape, text=black!65},
    wlab/.style={font=\footnotesize\itshape},
  ]
  % --- input layer ---
  \node[xnode] (i1) at (0, 2.0) {$x_1$};
  \node[xnode] (i2) at (0, 1.0) {$x_2$};
  \node[xnode] (i3) at (0, 0.0) {$x_3$};
  \node[dots]  (id) at (0,-1.0) {$\vdots$};
  \node[xnode] (in) at (0,-2.0) {$x_d$};

  % --- k teacher directions (k < d; linear projections, each its own color) ---
  \node[dnode=nethid]  (d1) at (2.4, 0.85) {};
  \node[dots]          (dd) at (2.4, 0.0)  {$\vdots$};
  \node[dnode=nethidc] (dk) at (2.4,-0.85) {};

  % --- general nonlinear link ---
  \node[hbox] (h) at (4.5, 0) {$h^\star$};

  % --- first-layer (W^*) edges, colored per direction ---
  \foreach \b/\c in {d1/nethid, dk/nethidc}
    \foreach \a in {i1,i2,i3,in}
      \draw[wedge=\c] (\a) -- (\b);

  % --- directions feed the link jointly ---
  \draw[aedge] (d1) -- (h.west);
  \draw[aedge] (dk) -- (h.west);

  % --- output ---
  \draw[oedge] (h.east) -- ++(1.4,0) node[right, font=\small, text=black] {$f^\star(\vec{x})\in\mathbb{R}$};

  % --- labels ---
  \node[wlab, text=black!70] at (1.2, 2.55) {$W^\star$};
  \node[lab] at (0,-3.0)   {$\vec{x}\sim\mathcal{N}(0,I_d)$};
  \node[lab] at (2.4,-3.0) {$k$ directions};
  \node[lab] at (4.5,-3.0) {link $h^\star$};
\end{tikzpicture}
}\par\smallskip
    {\footnotesize (a) multi-index teacher}
  \end{minipage}\hfill
  \begin{minipage}[b]{0.49\linewidth}
    \centering
    \resizebox{\linewidth}{!}{% Schematic of the two-layer (shallow) STUDENT network
%   f(x) = (1/p) sum_j a_j sigma(<w_j, x>)  -- panel (b) of the teacher-student
% figure in chapters/introduction.tex.  Standalone tikzpicture, resized to its
% minipage at the call site.  Needs tikz library arrows.meta (loaded in thesis.tex).
\begin{tikzpicture}[
    >={Latex[length=1.8mm]},
    font=\small,
    xnode/.style={circle, draw=netin!75, fill=netin!12, minimum size=7mm, inner sep=0pt},
    hnode/.style={circle, draw=#1!90, fill=#1!20, minimum size=8.5mm, inner sep=0pt},
    hnode/.default=nethid,
    onode/.style={circle, draw=netout!95, fill=netout!32, minimum size=9.5mm, inner sep=0pt},
    dots/.style={draw=none, fill=none, inner sep=1pt},
    wedge/.style={draw=#1!55, line width=0.5pt},
    wedge/.default=nethid,
    aedge/.style={draw=netout!75, line width=0.8pt},
    oedge/.style={->, draw=netout, line width=0.9pt},
    lab/.style={font=\footnotesize\itshape, text=black!65},
    wlab/.style={font=\footnotesize\itshape},
  ]
  % --- input layer ---
  \node[xnode] (i1) at (0, 2.0) {$x_1$};
  \node[xnode] (i2) at (0, 1.0) {$x_2$};
  \node[xnode] (i3) at (0, 0.0) {$x_3$};
  \node[dots]  (id) at (0,-1.0) {$\vdots$};
  \node[xnode] (in) at (0,-2.0) {$x_d$};

  % --- hidden layer (each neuron + its incoming weights share a color) ---
  \node[hnode=nethid]  (h1) at (2.6, 1.5) {$\sigma$};
  \node[hnode=nethidb] (h2) at (2.6, 0.5) {$\sigma$};
  \node[dots]          (hd) at (2.6,-0.5) {$\vdots$};
  \node[hnode=nethidc] (hp) at (2.6,-1.5) {$\sigma$};

  % --- output ---
  \node[onode] (o) at (5.2, 0) {$\textstyle\sum$};

  % --- first-layer (W) edges, colored per neuron ---
  \foreach \b/\c in {h1/nethid, h2/nethidb, hp/nethidc}
    \foreach \a in {i1,i2,i3,in}
      \draw[wedge=\c] (\a) -- (\b);

  % --- second-layer (a) edges ---
  \foreach \b in {h1,h2,hp}
    \draw[aedge] (\b) -- (o);

  % --- readout arrow ---
  \draw[oedge] (o) -- ++(1.5,0) node[right, font=\small, text=black] {$f(\vec{x})\in\mathbb{R}$};

  % --- labels ---
  \node[wlab, text=black!70] at (1.3, 2.55) {$W$};
  \node[wlab, text=netout]   at (4.0, 1.15) {$\vec{a}$};
  \node[lab] at (0,-3.0)   {$\vec{x}\in\mathbb{R}^{d}$};
  \node[lab] at (2.6,-3.0) {$p$ neurons};
  \node[lab] at (5.2,-3.0) {readout};
\end{tikzpicture}
}\par\smallskip
    {\footnotesize (b) two-layer student}
  \end{minipage}
  \caption[The teacher--student scenario.]{%
    The teacher--student scenario.
    \textbf{(a)} The multi-index teacher~\eqref{eq:intro:multi-index}: the Gaussian
    input $\vec{x}\sim\mathcal{N}(0,I_d)$ is projected onto the $k<d$ teacher
    directions $W^\star$ (each drawn in its own color) and passed jointly through
    a general nonlinear link $h^\star$ to produce the label.
    \textbf{(b)} The two-layer student~\eqref{eq:intro:student}: each of the $p$
    hidden neurons applies the \emph{pointwise} activation $\sigma$ to a linear
    readout $\langle\vec{w}_j, \vec{x}\rangle/\sqrt{d}$ through its first-layer weights
    $\vec{w}_j$ (a row of $W$, each neuron in its own color), and the
    second-layer weights $\vec{a}$ combine them with the mean-field
    normalization $\sfrac{1}{p}$.
    The student is the special case in which the teacher's link is a sum of
    pointwise activations.}
  \label{intro:fig:teacher-student}
\end{figure}

Two-layer multi-index networks are just one among infinitely many possible ways to proxy deep learning. As non-exhaustive examples, in the literature we find deep linear networks \cite{saxe2014exact, arora2018optimization}, infinite-depth-and-width limits with DMFT \cite{bordelon2024depthwise}, and tensor programs \cite{yang2019tensor}. The approach we discuss is not the only one for two-layer networks either: diagonal linear networks \cite{pesme2021implicit,even2023sgd}, the mean-field limit of two-layer nets \cite{mei2018mean,chizat2018global}, the neural tangent kernel / lazy regime \cite{jacot2018neural}, and the landscape geometry of linear networks \cite{trager2020pure}.

The $\sfrac1p$ prefactor is the \emph{mean-field normalization}, which keeps the output
finite as $p \to \infty$.
The student parameters are $\Theta = (\vec{a}, W)$.
Crucially, Equation~\eqref{eq:intro:student} has the same functional dependence on a low-dimensional projection of the input as the
multi-index teacher~\eqref{eq:intro:multi-index}: the student is a special case of
the teacher.
This symmetry is why the setting is called the \emph{teacher--student scenario}.
The two-layer network is the simplest architecture that generates a non-convex
loss landscape; already with $k = p = 1$ and a nonlinear activation, the landscape
has saddle points and spurious local minima that make the dynamics of gradient
descent non-trivial.
\addnote{Think of \emph{phase retrieval} ($\sigma(z) = z^2$): the loss landscape has at least two global minima, at $\vec{w} = \pm \vec{w}^\star$, and a saddle point at $\vec{w} = 0$.}

To measure the quality of the prediction of the model, we use the quadratic (\emph{square}) loss \[\ell(\hat{y}, y) = \tfrac{1}{2}(\hat{y} - y)^2\]
throughout, and measure the overall performance with \emph{generalization error} or
\emph{population risk},
\begin{equation}
\label{eq:intro:risk}
    \mathcal{R}(W, \vec{a})
    = \mathbb{E}_{(\vec{x},\, y)}\!\left[\frac{1}{2}
        \bigl(f(\vec{x};\, W, \vec{a}) - y\bigr)^2\right].
\end{equation}
Note that the expectation is taken over the data distribution, not the training set; this is what makes it a measure of generalization, since we know the precise data distribution.
The training happens on samples drawn from the data distribution, and not on the population risk, which is not accessible to the learning algorithm.

The square loss---predicting a scalar $y \in \mathbb{R}$---is especially convenient:
its gradient is proportional to the residual error, the population risk admits a
closed-form expression in a small set of order parameters (defined below), and the
resulting loss landscape, while non-convex, is structured enough to be studied
rigorously.

\paragraph{What algorithm do we use?}
Stochastic gradient descent (SGD) is the workhorse optimizer of modern deep
learning: virtually every model in practice is trained with it or one of its
descendants.
Our proxy for SGD is \emph{online (one-pass) SGD}: at each step $\nu$ a fresh,
independently drawn sample $(\vec{x}^\nu, y^\nu)$ is used to form a gradient
estimate and update the first-layer weights,
\begin{equation}
\label{eq:intro:sgd}
    \vec{w}_j^{\,\nu+1} = \vec{w}_j^{\,\nu}
    - \gamma\,\nabla_{\vec{w}_j}\,\ell\!\left(
        f(\vec{x}^\nu;\, W^\nu, \vec{a}),\, y^\nu\right),
    \qquad j \in [p],
\end{equation}
with learning rate $\gamma > 0$.
The number of steps $n$ equals the number of samples consumed.
Modern practice uses mini-batch multi-epoch training with sophisticated optimizers
(Adam, Muon, etc.) on datasets that are revisited many times.
Ours is simpler in the one key respect that makes it analytically tractable: the
\emph{independence} of successive gradient steps.
This independence is also not as artificial as it might seem: when the dataset is
large enough that each sample is seen at most once, the training regime is
effectively online.
\addnote{We relax the independence assumption in Chapter~\ref{chap:paper-2405}, and
lifting it turns out to be a pivotal step in the theory.}

Finally, we initialize the weights randomly: $\vec{w}_j^0 \sim \mathcal{N}(0, I_d)$.


\paragraph{Order parameters.}
We now come to the central observation that makes the whole program analytically
tractable: both teacher and student see the input only through a handful of
one-dimensional projections, the \emph{pre-activations} (or \emph{local fields})
\begin{equation}
\label{eq:intro:preactivations}
    \vec{\lambda} = \frac{W \vec{x}}{\sqrt{d}} \in \mathbb{R}^p, \qquad
    \vec{\lambda}^\star = \frac{W^\star \vec{x}}{\sqrt{d}} \in \mathbb{R}^k,
\end{equation}
where $W$ and $W^\star$ are the student and teacher weight matrices respectively
(defined above).
Indeed, the student~\eqref{eq:intro:student} depends on the input only through its
own pre-activations: writing $\lambda_j = \langle \vec{w}_j, \vec{x}\rangle/\sqrt{d}$, we have
\begin{equation}
\label{eq:intro:student-preactivations}
    f(\vec{x};\, W, \vec{a}) = h(\vec{\lambda}, \vec{a}), \qquad
    h(\vec{\lambda}, \vec{a}) = \frac{1}{p}\sum_{j=1}^p a_j\,\sigma(\lambda_j),
\end{equation}
which mirrors exactly the teacher $f^\star(\vec{x}) = h^\star(\vec{\lambda}^\star)$
from~\eqref{eq:intro:multi-index}.

Because the input is Gaussian and the pre-activations are linear in $\vec{x}$, the
pair $(\vec{\lambda}, \vec{\lambda}^\star)$ is \emph{jointly Gaussian}, with zero mean
and a covariance fixed entirely by the weight matrices, not by any individual data
point.
That covariance is assembled from the pairwise inner products of the student and
teacher rows, which we collect into the \emph{order parameters} (or \emph{overlap
matrices}):
\begin{equation}
\label{eq:intro:overlaps}
    Q_{jl} = \frac{\langle \vec{w}_j,\, \vec{w}_l\rangle}{d}, \qquad
    M_{jr} = \frac{\langle \vec{w}_j,\, \vec{w}^\star_r\rangle}{d}, \qquad
    P_{rs} = \frac{\langle \vec{w}^\star_r,\, \vec{w}^\star_s\rangle}{d}.
\end{equation}
The matrix $Q \in \mathbb{R}^{p \times p}$ encodes the self-correlations among
student neurons; $M \in \mathbb{R}^{p \times k}$ measures the alignment of each
student neuron with each teacher direction; and $P \in \mathbb{R}^{k \times k}$
describes the geometry of the teacher (concentrating on $P = I_k$ in the
high-dimensional limit under our Gaussian normalization of the teacher directions).
Together they form the covariance of the pre-activations,
\begin{equation}
\label{eq:intro:Omega}
    (\vec{\lambda}, \vec{\lambda}^\star) \sim \mathcal{N}(0, \Omega), \qquad
    \Omega = \begin{pmatrix} Q & M \\ M^\top & P \end{pmatrix}
    \in \mathbb{R}^{(p+k) \times (p+k)}.
\end{equation}
The triple $(Q, M, \vec{a})$---equivalently $(\Omega, \vec{a})$, since $P$ is
fixed---is therefore a set of \emph{sufficient statistics} for the problem: every
quantity that depends on the weights only through the joint law of the
pre-activations is a function of it alone.

The population risk is exactly such a quantity.
Substituting the student~\eqref{eq:intro:student-preactivations} into the
risk~\eqref{eq:intro:risk} and using the joint Gaussianity~\eqref{eq:intro:Omega},
the expectation over $\vec{x} \in \mathbb{R}^d$ collapses to an expectation over the
$(p+k)$-dimensional Gaussian $\mathcal{N}(0, \Omega)$:
\begin{equation}
\label{eq:intro:risk-preactivations}
    \mathcal{R}(\Omega, \vec{a})
    = \mathbb{E}_{(\vec{\lambda},\, \vec{\lambda}^\star) \sim \mathcal{N}(0, \Omega)}
        \!\left[\frac{1}{2}
        \bigl(h(\vec{\lambda}, \vec{a}) - h^\star(\vec{\lambda}^\star)\bigr)^2\right].
\end{equation}
The crucial feature of this expression is that it \emph{does not scale with $d$}: it
involves only the $O(p^2 + pk)$ entries of $(\Omega, \vec{a})$, with the ambient
dimension entering nowhere.
Rather than tracking the $p \times d$ weight vectors and an expectation over
$\mathbb{R}^d$, we track a handful of scalar overlaps and integrate a fixed
low-dimensional Gaussian.
It is precisely this dimension-independence that lets us send $d \to \infty$, the
core missing piece of our analysis.
